\documentclass{article}
\usepackage{amsmath}
\usepackage{amssymb}
\usepackage{geometry}
\usepackage{amsthm}
\usepackage{booktabs}
\usepackage{graphicx}  % Add this package
\usepackage{titling}

\geometry{a4paper, margin=1in}

\title{Foundational Mathematical Concepts for AI Futures}
\author{Dr James Carney}
\date{}

% Custom title format
\pretitle{\begin{center}\includegraphics[width=0.2\textwidth]{Full_logo_black.png}\\[0.5cm]\Huge\bfseries}
\posttitle{\end{center}}
\preauthor{\begin{center}\Large}
\postauthor{\end{center}}
\predate{\begin{center}\large}
\postdate{\end{center}}
\begin{document}

\maketitle


\vspace{1cm}

\section*{}



\subsection*{1. Equation of a Line}
The equation of a straight line in two-dimensional Cartesian coordinates is given by:

\[
y = mx + b
\]

where:
\begin{itemize}
    \item $y$ is the dependent variable (vertical axis value)
    \item $x$ is the independent variable (horizontal axis value)
    \item $m$ is the \textbf{slope} of the line, representing the rate of change
    \item $b$ is the \textbf{y-intercept}, the point where the line crosses the y-axis ($x = 0$)
\end{itemize}

The slope $m$ can be calculated from any two points $(x_1, y_1)$ and $(x_2, y_2)$ on the line:
\[
m = \frac{y_2 - y_1}{x_2 - x_1} \quad \text{for } x_1 \neq x_2
\]

\subsection*{2. Mean Squared Error (MSE)}
Given a set of $n$ data points with true values $y_i$ and predicted values $\hat{y}_i$, the Mean Squared Error measures the average squared difference between predictions and actual values:

\[
\text{MSE} = \frac{1}{n} \sum_{i=1}^{n} (y_i - \hat{y}_i)^2
\]

where:
\begin{itemize}
    \item $\sum$ (sigma) is the summation operator, meaning "add the results of this operation for every value of $i$ from 1 to $n$ for $y_i$ and $\hat{y}_i$"
    \item $n$ is the number of data points
    \item $y_i$ is the actual (true) value of the $i$-th observation
    \item $\hat{y}_i$ is the predicted value of the $i$-th observation
    \item $(y_i - \hat{y}_i)$ is the \textbf{residual} (prediction error)
\end{itemize}

MSE is always non-negative ($\text{MSE} \geq 0$), with smaller values indicating better predictions. Squaring the errors has two main effects:
\begin{enumerate}
    \item Ensures all differences are positive
    \item Gives more weight to larger errors (penalizes outliers)
\end{enumerate}

\subsection*{3. Power Rule of Differentiation}
The power rule is a fundamental technique for finding derivatives of polynomial functions. Given a function $f(x) = x^n$ where $n$ is a real number:

\[
\frac{d}{dx}(x^n) = n x^{n-1}
\]

\subsubsection*{Geometric Interpretation}
The derivative $\frac{df}{dx}$ represents the \textbf{slope of the tangent line} to the curve $f(x)$ at a specific point $x_0$:

\[
\text{Slope at } x_0 = f'(x_0) = \left. \frac{df}{dx} \right|_{x=x_0}
\]

This slope $f'(x_0)$ is exactly analogous to the coefficient $m$ in the line equation $y = mx + b$, but applied to the tangent line at a specific point on a curve.

\subsubsection*{Examples of the Power Rule}
\begin{align*}
    \frac{d}{dx}(x^2) &= 2x^{1} = 2x \\
    \frac{d}{dx}(x^3) &= 3x^{2} \\
    \frac{d}{dx}(x) &= 1\\
    \frac{d}{dx}(\text{constant}) &= 0 \quad \text{(since a constant number has no rate of change by defintition)}
\end{align*}

\subsubsection*{Visual example}
For a function $f(x) = x^3 - 6x^2 + 4x + 12$, the derivative is given by:
\[
\frac{d}{dx}f(x) = f'(x)= 3x^2 - 12x + 4 
\]

The derivative function $f'(x)$ gives the slope of the tangent at every point on $f(x)$. See Figure~\ref{fig:function_plot} for a visual illustration of this.

\begin{figure}[h!]  % h! = "here!" (try to place it here)
    \centering
    \includegraphics[width=0.8\textwidth]{tangents.jpg}
    \caption{Graph of $f(x) = x^3 - 6x^2 + 4x + 12$ and two tangents. The derivative  $f'(x)= 3x^2 - 12x + 4$ gives the slope of these tangents}
    \label{fig:function_plot}
\end{figure}

\section*{Summary of Relationships}
\begin{center}
\begin{tabular}{@{}lll@{}}
\toprule
\textbf{Concept} & \textbf{Mathematical Form} & \textbf{Interpretation} \\
\midrule
Line Equation & $y = mx + b$ & Describes a straight line \\
MSE & $\frac{1}{n}\sum (y_i - \hat{y}_i)^2$ & Measures prediction accuracy \\
Power Rule & $\frac{d}{dx}(x^n) = nx^{n-1}$ & Finds instantaneous rate of change \\
\bottomrule
\end{tabular}
\end{center}

\section*{Practical Connection}
In machine learning, particularly linear regression:
\begin{enumerate}
    \item We use the \textbf{line equation} $y = mx + b$ as our prediction model
    \item We measure prediction quality using \textbf{MSE} as our loss function
    \item We find optimal $m$ and $b$ values by minimizing MSE, often using derivatives obtained via the \textbf{power rule} (and other differentiation rules)
\end{enumerate}

\end{document}